\section{A bounded-rank Iwasawa module for the cubic field}
\label{sec:cubic-iwasawa-tower}

The fixed coefficient-prime-three sequence admits a finite-dimensional
arithmetic description.  This description does not estimate a
specialization by declaring its characteristic series nonzero.
Retain $E=\mathbb Q(\omega,\sqrt[3]{269})$, and set
\[
 K_n=\mathbb Q(\mu_{3^{n+1}}),\quad E_n=EK_n,\quad
 A_n=\operatorname{Cl}(E_n)_{(3)},\quad
 X=\varprojlim_n A_n,
\]
where transition maps are norms.  Let $A'_n$ be the quotient of $A_n$
by the classes of primes above three, and put
$X'=\varprojlim_n A'_n$.  The latter is the split unramified module
in (KL.9).  Write $\sigma$ for a generator of
$\operatorname{Gal}(E_n/K_n)$, chosen compatibly, and let
$\mathcal O=\mathbb Z_3[\omega]$ act by $\omega=\sigma$.

\begin{theorem}[AC-K16: seven generators suffice in the whole tower]
\label{AC-K16}
For every $n\geq0$,
\[
 \operatorname{Cl}(K_n)_{(3)}=0,\qquad
 \#\{\mathfrak p\subset\mathcal O_{K_n}:\mathfrak p\mid269\}
       =3^{\min(n,2)}.                                  \tag{KI.1}
\]
The norm $1+\sigma+\sigma^2$ annihilates $A_n$, and
\[
 \dim_{\mathbb F_3}(A_n/\lambda A_n)
 =
 \begin{cases}
 0,&n=0,\\
 2,&n=1,\\
 7,&n\geq2,
 \end{cases}
 \qquad \lambda=\omega-1.                              \tag{KI.2}
\]
All transition norms $A_{n+1}\to A_n$ are surjective.
Both $X$ and $X'$ are finitely generated $\mathcal O$-modules.
The minimal number of generators of $X$ is exactly seven,
and that of $X'$ is either six or seven.  In particular they are finitely
generated over $\mathbb Z_3$, of rank at most fourteen; their
Iwasawa $\mu$-invariants vanish.  No assumption of A or of a
main conjecture enters these assertions.
\end{theorem}
\begin{proof}
\proofstep{1. The cyclotomic base and the ramified primes.}
The class number of $K_0=\mathbb Q(\omega)$ is one.
For $n\geq1$, the extension $K_n/K_{n-1}$ is cyclic of degree
three, with just one ramified finite prime, the unique prime
above three, totally ramified.  There are no real places.
Suppose inductively that the three-primary class group of
$K_{n-1}$ is zero.  The three-primary part of Chevalley's
ambiguous class-number formula bounds the invariant three-class
group of $K_n$ by $3/3=1$, since its remaining unit index is
an integer at least one.  A nonzero finite abelian three-group
acted on by a group of order three has nonzero fixed points:
on its nonzero three-torsion vector space, $(\sigma-1)^3=0$.
Thus the entire three-class group of $K_n$ is zero.  This proves
the first assertion in (KI.1) by induction.  The exact known
formula used is C.-F.~Yu, \emph{Ambiguous class number formulas},
MPIM preprint 14--68 (2014), Theorem~1.1, equation~(1.2),
printed p.~1, with the ordinary-class interpretation on p.~2.
Its application here is the same ordinary, totally imaginary
case used in AC-K12.

Since $269\equiv-1\pmod{27}$ and
$v_3(269^2-1)=3$, lifting the exponent gives
\[
 \operatorname{ord}_{3^{n+1}}(269)
       =2\,3^{\max(n-2,0)}.                             \tag{KI.3}
\]
For completeness, an odd power of $269$ is $-1$ modulo three;
for an even exponent $2b$ its distance from one has valuation
$3+v_3(b)$.  The smallest exponent producing one modulo
$3^{n+1}$ is exactly (KI.3).  Cyclotomic Frobenius is exponentiation
by $269$, so dividing $[K_n:\mathbb Q]=2\,3^n$ by (KI.3)
proves the second assertion in (KI.1).

The intersection $E\cap K_\infty$ is $K_0$: the cubic extension
$E/K_0$ ramifies at $269$, whereas $K_\infty/K_0$ does not.
Thus every $E_n/K_n$ has degree three.  At each prime over
$269$, the valuation of $269$ is one, and $X^3-269$ is
Eisenstein; these are totally ramified primes of index three.
The Hensel calculation in AC-G22 puts all three roots in
$\mathbb Q_3(\omega)$, so every prime over three splits in
$E_n/K_n$.  The discriminant of the defining cubic shows there
is no further ramification.  Consequently $E_n/K_n$ has exactly
$t_n=3^{\min(n,2)}$ ramified primes, all tame and finite.

\proofstep{2. A uniform genus bound, with the root-of-unity saving.}
For an ideal class in $A_n$, extension after norm to $K_n$ is
the action of $1+\sigma+\sigma^2$.  The norm class has three-power
order and the base three-class group is zero, so this action
vanishes.  The action therefore factors through
\[
 \mathbb Z_3[C_3]/(1+\sigma+\sigma^2)\simeq\mathcal O.
\]
Chevalley's formula gives
\[
 \#A_n^{\langle\sigma\rangle}
 =\frac{3^{t_n-1}}
 {[\mathcal O_{K_n}^{\times}:
   \mathcal O_{K_n}^{\times}\cap
              N_{E_n/K_n}(E_n^{\times})]_{(3)}}.        \tag{KI.4}
\]
Here $[\cdot]_{(3)}$ denotes the three-primary part of the
finite index.  In fact this unit quotient has exponent dividing
three, since the cube of any base unit is a norm.
Without using the denominator, (KI.4) gives the bounds zero
and two on the logarithms for $n=0,1$.

For $n\geq2$ the denominator is at least three.  Indeed, choose
a primitive root $\zeta_{3^{n+1}}\in K_n$.  At a prime over
$269$ the residue-field size is
$q=269^{2\,3^{n-2}}$, and (KI.3) gives
$v_3(q-1)=n+1$.  Reduction preserves the exact order
$3^{n+1}$ of this root of unity, because that order is prime
to $269$.  Hence its residue is not a cube in the residue-field
unit group: the three-part of that group has precisely order
$3^{n+1}$.  In a totally ramified tame extension of degree
three, reduction of the norm of a unit is its residue cubed.
Also a norm which is a unit must come from a unit, since
$v_{K_n}(Nx)=f\,v_{E_n}(x)$ and $f=1$.
Thus $\zeta_{3^{n+1}}$ is not even a local norm, and so cannot
be a global norm.  The unit quotient in (KI.4) is nontrivial
and has order at least three.  This proves
$\#A_n^{\langle\sigma\rangle}\leq3^7$.

The invariant group is $A_n[\lambda]$.  For a finite
$\mathcal O$-module, multiplication by $\lambda$ has kernel
and cokernel of the same cardinality.  Both are vector spaces
over $\mathcal O/\lambda=\mathbb F_3$.
The upper bounds in (KI.2) now follow from (KI.4).

We next check that the unit index in (KI.4) is exactly three
for $n\geq2$ and exactly one for $n=0,1$.
For any unit $u$ in the CM field $K_n$, $u/\overline u$
is a root of unity.  Indeed it is an algebraic integer unit
all of whose conjugates have modulus one; the elementary
Kronecker lemma follows by bounding the integral coefficients
of the minimal polynomials of its powers, so that two powers
coincide.  Thus $u^2$ is a root of unity times the real unit
$u\overline u$.  In a quotient of exponent three, squaring
is bijective.  It is therefore enough to examine roots of
unity and real units.

At a prime over $269$, complex conjugation is the unique
order-two residue Frobenius: it is exponentiation by
$269^{3^{\max(n-2,0)}}$.  A real unit reduces to the
subfield of that cardinality.  Its multiplicative group has
order prime to three, since this odd power of $269$ is
$-1$ modulo three.  Hence every real-unit residue is a cube.
In this tame local extension that is sufficient for being
a unit norm: principal units are cubes, and lifting a
cubic residue reduces to this assertion.  At every other
finite place the extension is unramified or split, so units
are local norms.  There is no condition at a complex place.
The cyclic Hasse norm theorem now makes every real unit
a global norm.  Its exact locator is Milne, \emph{Class Field
Theory}, version~4.03, VIII, Theorem~3.1, printed p.~234.
For unramified local unit norms one may use V,
Proposition~4.12(c), printed p.~177.

The roots of unity of $K_n$ have order $2\,3^{n+1}$.
Their norm obstruction has order at most three, and the
primitive root just tested shows it has order three for
$n\geq2$.  At $n=0,1$, the residue field has cardinality
$269^2$ and its cube subgroup already contains all the
$3^{n+1}$-th roots.  These roots are local norms everywhere
and hence global norms as well.  Thus the unit index is
one in those two cases.  Substitution in (KI.4) proves
all the equalities in (KI.2).

\proofstep{3. Surjective norms and the compact limit.}
The extension $E_{n+1}/E_n$ is cyclic of degree three and
is totally ramified at each of its primes above three: its
local extension there is the cyclotomic step
$\mathbb Q_3(\mu_{3^{n+2}})/\mathbb Q_3(\mu_{3^{n+1}})$.
Let $H_n/E_n$ be the Hilbert three-class field.  Since $H_n/E_n$
is unramified, $H_n\cap E_{n+1}=E_n$.
Its base change $H_nE_{n+1}/E_{n+1}$ remains unramified and
abelian, and has Galois group $\operatorname{Gal}(H_n/E_n)$.
The Hilbert three-class field of $E_{n+1}$ consequently
surjects onto this group.  Under Artin reciprocity this
restriction is exactly the ideal-class norm; hence
$A_{n+1}\to A_n$ is onto.  The known reciprocity and
functoriality statements used are J.~S.~Milne,
\emph{Class Field Theory}, version~4.03 (2020), Chapter~V,
Proposition~3.3 and Theorems~3.5--3.6, printed pp.~157--158.
One can check the norm compatibility directly on an unramified
prime using the power of Frobenius in Proposition~3.3.

The maps $A_{n+1}/\lambda A_{n+1}\to A_n/\lambda A_n$
are therefore onto, and their dimensions form a nondecreasing
sequence bounded by seven.  They are exactly seven from level two onward.
There is a canonical isomorphism
\[
 X/\lambda X\simeq\varprojlim_n A_n/\lambda A_n.       \tag{KI.5}
\]
Here is a compactness verification to make the passage exact.
The inverse limit of any exact sequence of finite abelian
systems is exact: for a coherent element of the quotient,
its possible lifts for each finite initial segment form
nonempty finite sets, and the closed finite-constraint sets
in the product have the finite intersection property.
Similarly, if $y_n\in\lambda A_n$ is a coherent sequence,
solutions of $\lambda x_n=y_n$ over a finite initial segment
are nonempty (choose a solution at its largest level and
apply norms).  Compactness gives compatible solutions at all
levels.  Thus $\varprojlim\lambda A_n=\lambda X$, which
proves (KI.5).

Lift a basis of the finite vector space (KI.5) to at most
seven elements of $X$, and let $B$ be their $\mathcal O$-span.
It is closed, being a continuous image of a finite product
of the compact ring $\mathcal O$.  The compact pro-three
quotient $C=X/B$ satisfies $C=\lambda C$.  Since
$\lambda^2=-3\omega$, it also satisfies $C=3C$.
Every finite continuous quotient of a compact pro-three
abelian group is a finite three-group; multiplication by
three cannot be surjective on such a group unless it is
zero.  Hence all finite quotients of $C$ are zero, which
forces $C=0$.  This proves finite generation and the bound.
The finite-system compactness argument also makes
$X\to\varprojlim A'_n=X'$ surjective.  Its kernel is
cyclic over $\mathcal O$.  To see this, choose one compatible
sequence of primes $\mathfrak p_n$ above three.  Total
ramification in each tower step gives
$N[\mathfrak p_{n+1}]=[\mathfrak p_n]$.  Their three-primary
classes therefore define $p\in X$.  At every level the
three prime classes are its $\sigma$-translates; their sum
is the class of the principal extended cyclotomic prime,
so their span is $\mathcal O p_n$.  For any coherent sequence
in these spans, the equations $b p_n=y_n$ have a simultaneous
solution $b\in\mathcal O$ by the finite-intersection property
in the compact ring $\mathcal O$.  Consequently the kernel
is exactly $\mathcal O p$.  Since (KI.5) has dimension seven,
quotienting by this cyclic module leaves a reduction modulo
$\lambda$ of dimension six or seven.  This proves the
asserted minimal generator counts.

Finally a finitely generated $\mathbb Z_3$-module with a
continuous cyclotomic action is a torsion Iwasawa module
without a factor supported at the prime $(3)$; equivalently
its $\mu$-invariant is zero.  Directly, its quotient modulo
three is finite-dimensional over $\mathbb F_3$, whereas any
nonzero $\Lambda/(3^a)$ elementary factor has infinite-dimensional
reduction modulo three.  This proves the last assertion.
\end{proof}

\begin{corollary}[AC-K16: one matrix and one finite correction]
\label{AC-K16-matrix}
Let $T=X'_{\rm tors}$, with torsion taken over $\mathcal O$,
and write $X'/T\simeq\mathcal O^r$, where $0\leq r\leq7$.
Choose a cyclotomic generator $\gamma$ with
$\varepsilon_3(\gamma)=4$, and let
$G\in\operatorname{GL}_r(\mathcal O)$ be its action on this
free quotient.  For each $e\geq1$,
\[
 \begin{split}
 a_e=m_e-1
  &=\operatorname{length}_{\mathcal O}
          \bigl(T/(4^e\gamma-1)T\bigr)
      +v_\lambda\det(4^eG-I_r),\\
 0&\leq\operatorname{length}_{\mathcal O}
          \bigl(T/(4^e\gamma-1)T\bigr)
      \leq\operatorname{length}_{\mathcal O}T.
 \end{split}                                             \tag{KI.6}
\]
The determinant is nonzero.  For $r=0$ it is interpreted as
one, and the sequence $m_e$ is uniformly bounded.
For positive $r$ this is a reduction to the evaluations of
one degree-at-most-seven polynomial, with fixed finite error.
No lower bound for those nonzero evaluations is asserted.
\end{corollary}
\begin{proof}
The torsion submodule of a finitely generated module over
$\mathcal O$ is finite and stable under $\gamma$.
A generator with cyclotomic value four exists since the
cyclotomic image is $1+3\mathbb Z_3$ and four generates that
group.  Apply the endomorphism $b_e=4^e\gamma-1$ to
$0\to T\to X'\to\mathcal O^r\to0$.
By (KL.9) its cokernel on $X'$ is the finite group $W_e(E)$.
Its cokernel on the free quotient is therefore finite, so
$\det(4^eG-I_r)\ne0$ and its kernel on that free quotient
is zero.  The snake lemma consequently gives the exact sequence
\[
 0\longrightarrow T/b_eT\longrightarrow W_e(E)
 \longrightarrow\mathcal O^r/(4^eG-I_r)\mathcal O^r
 \longrightarrow0.                                    \tag{KI.7}
\]
Smith normal form computes the length of the last group as
the determinant valuation.  Taking lengths, using (KL.1),
proves (KI.6).  The asserted error bound is simply that a
quotient of the fixed finite module $T$ has no greater length.
\end{proof}

\paragraph{Geometric meaning and the remaining estimate.}
At finite level, $A_n$ describes the maximal abelian
three-primary finite \'etale covering of
$\operatorname{Spec}\mathcal O_{E_n}$; quotienting by the
three-prime classes requires those primes to split.
This is the Galois-category setting of SGA~1, V~\S4 and
\S7, with arithmetic class field theory specifying its
abelian quotient.  The inverse limit is a compact
$\mathcal O$-module, not a finite-dimensional pure motive
or an arithmetic surface.  The Tate-twist specialization
is the wild-kernel comparison of (KL.9), related to higher
K-theory by AC-K10.  No Grothendieck archive scan asserts
the new generator bound or supplies the missing size estimate.
The root-of-unity norm argument is an integral ramification
calculation, not a positivity statement about the cubic
Weil distribution.  H04.4 and the comparison with Alp\"oge's
growing hosts remain open after (KI.6).
